\section{Implementation Details}
\label{sec:impl}

This section describes the software architecture, graphical user interface, technology stack, and critical code highlights, authored collaboratively by Members 1, 2, and 3.

\subsection{Software Architecture and Decoupled Subsystems}
The project follows a decoupled, object-oriented design adhering to SOLID principles. The system is partitioned into five distinct Python packages under \texttt{core/interfaces.py}:
\begin{enumerate}
    \item \texttt{env/}: Houses the Pygame environment, continuous physics simulation, collision engine, and procedural tilemap loaders for Levels 1--5.
    \item \texttt{discretizer/}: Implements \texttt{BaseDiscretizer}, encapsulating continuous coordinate projection into the discrete 7-tuple hashable representation ($\le 1{,}152$ states).
    \item \texttt{rewards/}: Implements the Strategy Pattern via \texttt{BaseRewardStrategy}, allowing plug-and-play swapping between Sparse, Dense, and Balanced reward formulations without touching environment physics.
    \item \texttt{agent/}: Implements \texttt{BaseAgent}, providing both \texttt{TabularQAgent} (pure off-policy Q-learning) and \texttt{HybridQSARSAAgent} (Expected SARSA, linear $Q(\sigma)$ blending, and hazard-gated switching), alongside experience rehearsal, Q-table serialization, and decay schedulers.
    \item \texttt{evaluation/}: Houses telemetry loggers, CSV export hooks, and automated Matplotlib/Seaborn graph generation scripts.
\end{enumerate}

\subsection{User Interface and Interactive Modes}
The application provides two execution modes:
\begin{itemize}
    \item \textbf{Interactive GUI / Evaluation Mode}: Renders at $60\text{ FPS}$ using Pygame. The head-up display (HUD) visualizes real-time metrics, including current episode index, instantaneous reward, exploration rate $\epsilon$, current discrete state tuple $\langle s_1, \dots, s_7 \rangle$, and selected action. A manual keyboard mode (\texttt{play\_human.py}) allows human benchmark play.
    \item \textbf{Headless Training Mode}: Bypasses graphical rendering to execute physics steps at over $500\text{ FPS}$, accelerating multi-episode curriculum completion.
\end{itemize}

\subsection{Tech Stack}
The system is implemented in Python 3.10 and relies on standard open-source libraries:
\begin{itemize}
    \item \textbf{Pygame 2.5}: Manages 2D sprite rendering, window events, delta-time clocks, and collision bounding boxes.
    \item \textbf{NumPy 1.26}: Handles matrix operations, Q-table indexing, candidate tie-breaking, and numerical array processing.
    \item \textbf{Pandas 2.2}: Manages tabular telemetry logging, data frame transformations, and rolling performance window calculations.
    \item \textbf{Matplotlib 3.8 \& Seaborn 0.13}: Renders publication-grade vector graphics (PDF) adhering to IEEE/Springer typographical styles.
\end{itemize}

\subsection{Code Highlights}
Listing~\ref{lst:tie_break} illustrates the deterministic uniform tie-breaking implementation in \texttt{agent/tabular\_q\_agent.py}, which eliminates early-stage action selection bias.

\begin{figure}[h]
\begin{center}
\begin{minipage}{0.9\textwidth}
\begin{verbatim}
def select_action(self, state: Tuple[int, ...]) -> int:
    if self.training and np.random.rand() < self.epsilon:
        return int(np.random.choice(self.num_actions))
    
    q_values = self.q_table[state]
    max_val = np.max(q_values)
    candidates = np.flatnonzero(q_values == max_val)
    return int(np.random.choice(candidates))
\end{verbatim}
\end{minipage}
\end{center}
\caption{Deterministic random tie-breaking during action selection.}
\label{lst:tie_break}
\end{figure}
